\documentclass[12pt,a4paper]{article}
\usepackage[utf8]{inputenc}
\usepackage[french]{babel}
\usepackage{amsmath}
\usepackage{amsfonts}
\usepackage{amssymb}
\usepackage{graphicx}
\usepackage{geometry}
\usepackage{fancyhdr}
\usepackage{hyperref}
\usepackage{xcolor}

\geometry{left=2.5cm, right=2.5cm, top=2.5cm, bottom=2.5cm}

\pagestyle{fancy}
\fancyhf{}
\lhead{Rapport RGPD}
\rhead{Protection des Données}
\rfoot{Page \thepage}

\hypersetup{
    colorlinks=true,
    linkcolor=blue,
    urlcolor=cyan,
}

\begin{document}

\begin{titlepage}
    \centering
    
    % En-tête officiel
    {\large \textbf{Ministère de l'Enseignement}}\\[0.2cm]
    {\large \textbf{Supérieur de la Recherche}}\\[0.2cm]
    {\large \textbf{Scientifique et de l'Innovation}}\\[0.5cm]
    
    {\large \textbf{BURKINA FASO}}\\[0.2cm]
    {\large \textbf{Unité - Progrès - Justice}}\\[1.5cm]
    
    % Logo (si disponible, sinon remplacer par le nom)
    {\Large \textbf{Université Joseph KI-ZERBO}}\\[0.5cm]
    {\large Faculté de Droit et Sciences Politiques}\\[1.5cm]
    
    % Titre principal
    {\Huge \textbf{RAPPORT DE PROJET}}\\[0.5cm]
    
    % Ligne verte supérieure
    {\color{green!60!black}\rule{0.8\textwidth}{3pt}}\\[0.3cm]
    {\LARGE \textbf{Règlement Général sur la Protection des Données}}\\[0.2cm]
    {\Large \textbf{(RGPD)}}\\[0.3cm]
    % Ligne verte inférieure
    {\color{green!60!black}\rule{0.8\textwidth}{3pt}}\\[2cm]
    
    % Membres du groupe
    \begin{flushleft}
    \hspace{2cm}
    {\large \textbf{Membres du Groupe :}}\\[0.5cm]
    \hspace{3cm} KABORE Wend-Benindo François\\[0.2cm]
    \hspace{3cm} SISSAO Sarata\\[0.2cm]
    \hspace{3cm} YÉ ELISÉ NIKIESSAN\\
    \end{flushleft}
    
    \vfill
    
    % Professeur
    \begin{flushright}
    \hspace{8cm}
    {\large \textbf{Professeur :}}\\[0.3cm]
    \hspace{8cm} Dr KOURAOGO\\
    \end{flushright}
    
    \vspace{1cm}
    
    % Année académique
    {\large \textbf{Année Académique : 2025 - 2026}}
    
\end{titlepage}

\tableofcontents
\newpage

\section*{Introduction}
\addcontentsline{toc}{section}{Introduction}

Avec l'essor du numérique et la multiplication des traitements automatisés, la question de la protection des données personnelles est devenue centrale. Les citoyens européens étaient confrontés à des pratiques divergentes selon les pays, ce qui créait une insécurité juridique et limitait la confiance dans les services numériques. Face à ces enjeux, l'Union européenne a adopté en 2016 le Règlement Général sur la Protection des Données (RGPD), afin d'instaurer un cadre unique et cohérent pour l'ensemble des États membres. Ce règlement constitue une avancée majeure dans la régulation du traitement des données à caractère personnel et dans la garantie des droits fondamentaux des individus.

La problématique essentielle posée par le RGPD est de savoir comment concilier la protection des droits fondamentaux des personnes physiques avec la nécessité de permettre la libre circulation des données dans un espace économique globalisé. Le texte répond à cette question en fixant des principes clairs et des obligations précises pour les acteurs du traitement, qu'il s'agisse des responsables ou des sous‑traitants, tout en prévoyant des mécanismes de contrôle et de sanction.

Les objectifs poursuivis par le règlement sont multiples. Il s'agit d'abord de garantir aux individus un contrôle renforcé sur leurs données personnelles, en leur donnant des droits étendus tels que le droit d'accès, de rectification, d'effacement ou encore de portabilité. Ensuite, le RGPD impose aux responsables de traitement et à leurs sous‑traitants des obligations strictes en matière de transparence, de sécurité et de responsabilité, afin d'assurer une gestion conforme et respectueuse des données. Enfin, il vise à favoriser la confiance dans l'économie numérique en assurant une protection uniforme dans toute l'Union européenne, créant ainsi un environnement propice au développement des échanges et des innovations technologiques.

\newpage

\section*{Chapitre I :Dispositions générales}
\addcontentsline{toc}{section}{Chapitre I :Dispositions générales}

Ce chapitre (Articles 1 à 4) établit le champ d'application du RGPD et ses objectifs. Il affirme que la protection des données personnelles est un droit fondamental et que leur libre circulation au sein de l'Union européenne ne doit pas être entravée. Le règlement s'applique à tout traitement de données, automatisé ou non, dès lors qu'il concerne des personnes physiques identifiées ou identifiables. Certaines exceptions sont prévues, notamment pour les activités strictement personnelles ou pour les traitements liés à la sécurité publique et aux enquêtes pénales.

Il introduit également les définitions clés : « données à caractère personnel », « traitement », « responsable du traitement », « sous‑traitant », « consentement », « profilage », « pseudonymisation », etc. Ces notions structurent l'ensemble du texte et clarifient les rôles et responsabilités des acteurs impliqués. Enfin, le champ territorial est élargi : le règlement s'applique aussi aux organisations situées hors UE dès lors qu'elles ciblent des personnes résidant sur le territoire européen.

En résumé, ce chapitre pose les bases du RGPD en définissant son périmètre, ses objectifs et son vocabulaire juridique, garantissant ainsi une compréhension commune et une application uniforme dans tous les États membres.

\newpage

\section*{Chapitre II :Principes}
\addcontentsline{toc}{section}{Chapitre II :Principes}

Le deuxième chapitre du RGPD fixe les principes fondamentaux qui doivent guider tout traitement de données personnelles. Ces principes constituent le cœur du règlement et définissent les conditions de licéité et de transparence.

Tout d'abord, les données doivent être traitées de manière licite, loyale et transparente vis‑à‑vis des personnes concernées. Elles doivent être collectées pour des finalités déterminées, explicites et légitimes, sans pouvoir être utilisées ultérieurement pour des objectifs incompatibles. Le règlement insiste sur la minimisation des données, c'est‑à‑dire que seules les informations strictement nécessaires doivent être collectées et traitées.

Un autre principe clé est l'exactitude : les données doivent être tenues à jour et corrigées en cas d'erreur. Elles ne peuvent être conservées que pour une durée limitée, proportionnée aux finalités du traitement, sauf exceptions pour des motifs d'intérêt public, de recherche scientifique ou statistique. Enfin, la sécurité est au centre du dispositif : les données doivent être protégées contre tout accès non autorisé, perte ou destruction, grâce à des mesures techniques et organisationnelles appropriées.

Le chapitre précise également les bases légales du traitement : consentement explicite de la personne concernée, exécution d'un contrat, respect d'une obligation légale, sauvegarde des intérêts vitaux, mission d'intérêt public ou poursuite d'intérêts légitimes. Le consentement doit être libre, éclairé et révocable à tout moment. Des dispositions spécifiques encadrent le consentement des enfants, notamment pour les services en ligne, avec un âge minimum fixé à 16 ans (ou 13 ans selon les États membres).

Enfin, le règlement interdit en principe le traitement des données sensibles (origine ethnique, opinions politiques, convictions religieuses, données biométriques ou de santé, orientation sexuelle), sauf exceptions strictement encadrées, comme le consentement explicite, la sauvegarde d'intérêts vitaux ou des motifs d'intérêt public majeur.

En résumé, ce chapitre établit les règles de base qui garantissent que tout traitement de données personnelles respecte la dignité, la liberté et les droits fondamentaux des individus, tout en permettant une utilisation proportionnée et sécurisée des informations.

\newpage

\section*{Chapitre III :Droits de la personne concernée}
\addcontentsline{toc}{section}{Chapitre III : Droits de la personne concernée}

Le troisième chapitre du RGPD est consacré aux droits accordés aux individus dont les données personnelles sont traitées. Il constitue l'un des piliers du règlement, car il renforce le pouvoir des citoyens sur leurs informations.

Les personnes concernées disposent d'un droit d'information : elles doivent être informées de manière claire et transparente sur la collecte et l'utilisation de leurs données. Elles ont également un droit d'accès, leur permettant de savoir quelles données sont détenues par un responsable de traitement et dans quel but. Le règlement prévoit aussi le droit de rectification, qui permet de corriger des données inexactes ou incomplètes.

Un droit majeur introduit par le RGPD est le droit à l'effacement, aussi appelé « droit à l'oubli ». Il permet à un individu de demander la suppression de ses données lorsque celles‑ci ne sont plus nécessaires, lorsqu'il retire son consentement ou lorsqu'elles ont été traitées de manière illicite. À cela s'ajoute le droit à la limitation du traitement, qui permet de suspendre temporairement l'utilisation des données dans certaines situations.

Le règlement instaure également le droit à la portabilité, une innovation importante : les individus peuvent obtenir leurs données dans un format structuré et lisible, afin de les transmettre facilement à un autre prestataire. Enfin, le droit d'opposition autorise une personne à refuser le traitement de ses données, notamment à des fins de prospection commerciale ou de profilage.

Ces droits sont complétés par des garanties spécifiques, comme le droit de ne pas faire l'objet d'une décision fondée uniquement sur un traitement automatisé, y compris le profilage, lorsqu'elle produit des effets juridiques significatifs.

En résumé, ce chapitre place l'individu au centre du dispositif en lui donnant des moyens concrets de contrôler ses données personnelles, renforçant ainsi la transparence et la confiance dans l'économie numérique.

\newpage

\section*{Chapitre IV : Responsable du traitement et sous traitant}
\addcontentsline{toc}{section}{Chapitre IV : Responsable du traitement et sous‑traitant}

Le quatrième chapitre du RGPD définit les obligations des acteurs principaux du traitement des données : le responsable du traitement et le sous‑traitant. Le responsable du traitement est celui qui détermine les finalités et les moyens du traitement, tandis que le sous‑traitant agit pour le compte du responsable. Ce chapitre établit un cadre de responsabilité renforcé afin d'assurer une gestion conforme et sécurisée des données personnelles.

Le responsable du traitement doit mettre en œuvre des mesures techniques et organisationnelles appropriées pour garantir la conformité au règlement. Cela inclut la tenue d'un registre des activités de traitement, la réalisation d'analyses d'impact sur la protection des données lorsque les risques sont élevés, et la mise en place de mécanismes de sécurité adaptés. Il doit également démontrer sa conformité en appliquant le principe de responsabilité proactive (« accountability »).

Le sous‑traitant, de son côté, ne peut agir que sur instruction du responsable et doit offrir des garanties suffisantes en matière de sécurité et de respect des droits des personnes concernées. Les relations entre responsables et sous‑traitants doivent être encadrées par un contrat écrit, précisant les obligations de chacun, notamment en matière de confidentialité, de sécurité et de coopération avec l'autorité de contrôle.

Ce chapitre introduit aussi la notion de délégué à la protection des données (DPO), obligatoire dans certains cas (par exemple pour les organismes publics ou les entreprises traitant des données sensibles à grande échelle). Le DPO a pour mission de conseiller, contrôler et servir de point de contact avec l'autorité de contrôle.

En résumé, le chapitre IV renforce la responsabilité des acteurs du traitement en imposant des obligations précises et en instaurant des mécanismes de contrôle. Il vise à garantir que la protection des données personnelles soit intégrée dès la conception des systèmes et tout au long de leur utilisation.

\newpage

\section*{Chapitre V : Transferts de données à caractère personnel vers des pays tiers ou à des organisations internationales}
\addcontentsline{toc}{section}{Chapitre V : Transferts de données vers des pays tiers}

Le cinquième chapitre du RGPD encadre les conditions dans lesquelles les données personnelles peuvent être transférées en dehors de l'Union européenne. L'objectif est de garantir que la protection accordée aux individus ne soit pas affaiblie lorsque leurs données quittent le territoire européen.

Le principe de base est que les transferts ne sont autorisés que si le pays tiers ou l'organisation internationale offre un niveau de protection adéquat, reconnu par une décision de la Commission européenne. Cette décision d'adéquation assure que les règles locales garantissent des droits et des garanties équivalents à ceux prévus par le RGPD.

En l'absence d'une telle décision, les transferts peuvent être réalisés sous réserve de garanties appropriées, comme des clauses contractuelles types, des règles d'entreprise contraignantes ou des accords spécifiques entre autorités publiques. Ces mécanismes visent à encadrer juridiquement le transfert et à protéger les droits des personnes concernées.

Le règlement prévoit également des dérogations limitées, par exemple lorsque la personne concernée a donné son consentement explicite, lorsque le transfert est nécessaire à l'exécution d'un contrat, à la sauvegarde d'intérêts vitaux ou encore à la constatation ou défense d'un droit en justice. Toutefois, ces dérogations doivent rester exceptionnelles et ne peuvent pas constituer une pratique courante.

En résumé, ce chapitre établit un équilibre entre la nécessité de permettre les échanges internationaux de données et l'exigence de maintenir un haut niveau de protection des droits fondamentaux. Il impose aux responsables de traitement une vigilance accrue et des obligations strictes lorsqu'ils envisagent de transférer des données hors de l'Union européenne.

\newpage

\section*{Chapitre VI :Autorités de contrôle indépendantes}
\addcontentsline{toc}{section}{Chapitre VI :Autorités de contrôle indépendantes}

Le sixième chapitre du RGPD définit le rôle et les missions des autorités de contrôle, qui sont des organismes publics indépendants créés par chaque État membre. Leur mission principale est de veiller à l'application correcte du règlement et de protéger les droits fondamentaux des personnes en matière de données personnelles.

Chaque autorité doit agir en toute indépendance, sans recevoir d'instructions extérieures, et disposer de moyens suffisants pour exercer ses fonctions. Le règlement précise que ces autorités doivent être dotées de pouvoirs étendus : elles peuvent enquêter, ordonner la rectification ou l'effacement de données, limiter ou suspendre un traitement, et infliger des sanctions administratives. Elles ont également un rôle de conseil auprès des institutions publiques et privées, et doivent sensibiliser les citoyens à leurs droits.

Le texte prévoit aussi la coopération entre autorités nationales, notamment dans les cas de traitements transfrontaliers. Pour éviter les divergences, un mécanisme de cohérence est instauré : les autorités doivent échanger des informations et travailler ensemble au sein du Comité européen de la protection des données (CEPD). Ce comité joue un rôle central en harmonisant l'interprétation du règlement et en émettant des lignes directrices.

En résumé, ce chapitre établit un système de contrôle indépendant et coordonné, garantissant que le RGPD soit appliqué de manière uniforme dans toute l'Union européenne. Les autorités de contrôle deviennent ainsi les garantes de la transparence, de la responsabilité et de la protection effective des droits des personnes.

\newpage

\section*{Chapitre VII :Coopération et cohérence}
\addcontentsline{toc}{section}{Chapitre VII :Coopération et cohérence}

Le septième chapitre du RGPD met en place un système destiné à assurer une application uniforme du règlement dans toute l'Union européenne. Il prévoit que les autorités de contrôle nationales doivent coopérer entre elles, notamment lorsqu'un traitement de données concerne plusieurs États membres. Cette coopération inclut l'échange d'informations, l'assistance mutuelle et la participation à des opérations conjointes.

Le texte introduit le mécanisme de cohérence, qui permet d'éviter des divergences d'interprétation. Lorsqu'une autorité nationale doit prendre une décision importante concernant un traitement transfrontalier, elle doit consulter les autres autorités concernées. En cas de désaccord, le Comité européen de la protection des données (CEPD) intervient pour rendre un avis ou une décision contraignante.

Ce comité, composé des représentants des autorités nationales et de la Commission européenne, joue un rôle central : il émet des lignes directrices, recommandations et bonnes pratiques afin d'harmoniser l'application du règlement. Il peut également arbitrer les différends entre autorités nationales.

En résumé, ce chapitre garantit que le RGPD soit appliqué de manière cohérente dans toute l'Union européenne, en instaurant un système de coopération obligatoire et un mécanisme de cohérence piloté par le CEPD.

\newpage

\section*{Chapitre VIII :Voies de recours, responsabilité et sanctions}
\addcontentsline{toc}{section}{Chapitre VIII :Voies de recours, responsabilité et sanctions}

Le huitième chapitre du RGPD établit les mécanismes permettant aux personnes concernées de défendre leurs droits et fixe les règles de responsabilité des acteurs du traitement. Il prévoit d'abord que tout individu qui estime que ses données ont été traitées en violation du règlement dispose d'un droit de recours effectif auprès d'une autorité de contrôle, d'un tribunal ou contre un responsable ou un sous‑traitant. Ce droit est garanti même si le traitement est effectué par une autorité publique.

Le texte précise ensuite la responsabilité civile des responsables de traitement et des sous‑traitants : toute personne ayant subi un dommage matériel ou moral du fait d'une violation du RGPD peut obtenir réparation. Le responsable du traitement est en principe tenu de réparer le préjudice, mais le sous‑traitant peut également être tenu responsable s'il n'a pas respecté ses obligations contractuelles ou réglementaires.

Ce chapitre introduit aussi le régime des sanctions administratives. Les autorités de contrôle disposent du pouvoir d'imposer des amendes proportionnées et dissuasives. Celles‑ci peuvent atteindre jusqu'à 20 millions d'euros ou 4 \% du chiffre d'affaires annuel mondial de l'entreprise, selon le montant le plus élevé. Le règlement distingue deux niveaux de sanctions : un premier plafond pour les manquements aux obligations générales (par exemple, absence de registre ou défaut de sécurité), et un second, plus élevé, pour les violations graves (comme le non‑respect des principes de base du traitement ou des droits des personnes).

En résumé, ce chapitre renforce la protection des individus en leur donnant des moyens concrets de recours et en imposant aux acteurs économiques un régime de responsabilité et de sanctions particulièrement strict, garantissant l'effectivité du RGPD.

\newpage

\section*{Chapitre IX :Dispositions relatives à des situations particulières de traitement}
\addcontentsline{toc}{section}{Chapitre IX :Dispositions relatives à des situations particulières}

Le neuvième chapitre du RGPD prévoit des règles spécifiques pour certaines situations particulières de traitement de données. Il reconnaît que, dans certains domaines, des ajustements sont nécessaires afin de concilier la protection des droits fondamentaux avec des impératifs d'intérêt public ou des besoins sectoriels.

Le texte autorise les États membres à introduire des dispositions particulières pour encadrer le traitement des données dans des contextes tels que la liberté d'expression et d'information, l'accès aux documents officiels, l'emploi et la sécurité sociale, la santé publique ou encore la recherche scientifique et statistique. Par exemple, des exceptions peuvent être prévues pour permettre aux journalistes de traiter des données personnelles dans le cadre de leur activité, tout en respectant l'équilibre entre vie privée et liberté d'information.

Ce chapitre aborde également les traitements liés à des missions d'intérêt public, comme la protection sociale, la sécurité nationale ou la justice. Dans ces cas, les États membres peuvent adapter certaines obligations du RGPD, à condition que ces dérogations soient proportionnées et respectent l'essence des droits fondamentaux.

En résumé, le chapitre IX introduit une flexibilité encadrée : il permet aux États membres de prévoir des règles spécifiques pour des secteurs sensibles ou stratégiques, tout en maintenant le principe général de protection des données personnelles.

\newpage

\section*{Chapitre X :Actes délégués et actes d'exécution}
\addcontentsline{toc}{section}{Chapitre X :Actes délégués et actes d'exécution}

Le dixième chapitre du RGPD traite des mécanismes juridiques permettant à la Commission européenne d'adopter des actes délégués et des actes d'exécution afin de préciser ou compléter certaines dispositions du règlement. Ces instruments garantissent que le texte puisse être adapté et appliqué de manière uniforme dans l'ensemble des États membres.

Les actes délégués permettent à la Commission de modifier ou compléter certains éléments non essentiels du règlement, par exemple pour préciser des critères techniques ou des modalités pratiques. Ils sont soumis à un contrôle du Parlement européen et du Conseil, qui peuvent révoquer ou s'opposer à ces actes si nécessaire.

Les actes d'exécution, quant à eux, visent à assurer une application cohérente du règlement dans toute l'Union. Ils permettent de définir des procédures communes ou des formulaires standards, afin que les autorités nationales appliquent les règles de manière harmonisée.

Ce chapitre illustre la volonté de l'Union européenne de maintenir une flexibilité encadrée : le RGPD reste un texte stable, mais il peut être ajusté par la Commission pour répondre aux évolutions technologiques ou aux besoins pratiques, tout en garantissant un contrôle démocratique par les institutions européennes.

En résumé, le chapitre X établit les outils juridiques nécessaires pour adapter et uniformiser l'application du RGPD, en combinant efficacité administrative et contrôle institutionnel.

\newpage

\section*{Chapitre XI :Dispositions finales}
\addcontentsline{toc}{section}{Chapitre XI :Dispositions finales}

Le onzième et dernier chapitre du RGPD regroupe les dispositions finales qui assurent la cohérence juridique et l'entrée en vigueur du règlement. Il précise d'abord que le RGPD abroge la directive 95/46/CE, qui constituait jusqu'alors le cadre européen de protection des données. Cette abrogation marque une rupture importante : le règlement, contrairement à une directive, est directement applicable dans tous les États membres, sans nécessiter de transposition nationale.

Le texte rappelle également que le règlement est d'intérêt pour l'Espace économique européen (EEE), ce qui signifie qu'il s'applique non seulement aux pays de l'Union européenne, mais aussi aux États de l'EEE, garantissant une harmonisation plus large.

Enfin, ce chapitre fixe la date d'entrée en application : le RGPD a été adopté le 27 avril 2016 et est entré en vigueur le 25 mai 2018, après une période de transition permettant aux entreprises et administrations de se préparer. Il souligne que toutes les dispositions doivent être interprétées en cohérence avec les droits fondamentaux consacrés par la Charte des droits fondamentaux de l'Union européenne.

En résumé, le chapitre XI clôt le texte en confirmant son caractère obligatoire et uniforme, en élargissant son champ d'application à l'EEE et en fixant son entrée en vigueur. Il consacre ainsi le RGPD comme le cadre juridique de référence en matière de protection des données personnelles en Europe.

\newpage

\section*{Conclusion}
\addcontentsline{toc}{section}{Conclusion}

Le Règlement Général sur la Protection des Données (RGPD) constitue une avancée majeure dans la protection des droits fondamentaux des citoyens européens. À travers ses onze chapitres, il établit un cadre juridique complet qui encadre le traitement des données personnelles, fixe des principes clairs de licéité, transparence et sécurité, et impose des obligations strictes aux responsables de traitement comme aux sous‑traitants.

Notre analyse a montré que le RGPD repose sur trois piliers essentiels :

La protection des individus, grâce à des droits renforcés tels que l'accès, la rectification, l'effacement ou la portabilité.

La responsabilité des acteurs, avec des obligations précises, la désignation de délégués à la protection des données et des sanctions dissuasives pouvant atteindre 20 millions d'euros ou 4 \% du chiffre d'affaires mondial.

La cohérence européenne, assurée par la coopération des autorités de contrôle et le rôle central du Comité européen de la protection des données.

Au‑delà des principes généraux, le règlement prévoit des dispositions adaptées à des situations particulières (liberté d'expression, santé publique, recherche scientifique), ainsi que des mécanismes de flexibilité encadrée via les actes délégués et d'exécution. Enfin, ses dispositions finales consacrent son caractère obligatoire et uniforme, en élargissant son champ d'application à l'Espace économique européen.

\end{document}
